\clearpage
\section{CONTEXTUALIZAÇÃO DA INSTITUIÇÃO}
Nos primeiros vinte anos após a Proclamação da República, as indústrias brasileiras já apresentavam algum crescimento, demandando a necessidade de mão de obra mais qualificada. As novas tarefas exigiam pessoas com conhecimentos especializados e apontavam para a necessidade de se estabelecer, de imediato, o ensino profissional.

Assim, em setembro de 1909, o então Presidente do Brasil, Nilo Peçanha, mediante Decreto-Lei n° 7.566, cria nas capitais dos estados da república, as Escolas de Aprendizes Artífices para o ensino profissional primário e gratuito.

A Instituição, batizada com a denominação de Escola de Aprendizes Artífices do Ceará, foi instalada no dia 24 de maio de 1910, na Av. Alberto Nepomuceno, onde funciona, atualmente, a Secretaria Estadual da Fazenda.

Em 1930 o governo provisório assume o poder e a educação passa a ser regulada pelo Ministério da Educação e Saúde Pública (MESP). As Escolas de Aprendizes Artífices, anteriormente ligadas ao Ministério da Agricultura, passaram, por consequência e de imediato, ao MESP e a receber subsídios do governo central.

Em 1937, na reforma do Ministério da Educação e Saúde Pública, o ministro Capanema, mediante a Lei n° 378 de 13 de janeiro, transforma as Escolas de Aprendizes Artífices em Liceus Profissionais, recebendo, no Ceará, a denominação de Liceu Industrial de Fortaleza.

Com a eclosão da Segunda Guerra Mundial, em primeiro de setembro de 1939, houve intensa redução na importação de produtos estrangeiros. Por esta razão, o Brasil passou a cuidar da implantação de indústrias básicas, incentivando a criação de estabelecimentos fabris e, consequentemente, adotou uma política paralela de incentivo à formação de mão-de-obra qualificada, para atender ao incipiente parque industrial.

Por despacho do Ministro da Educação, em 28 de agosto de 1941, houve uma outra modificação no nome dos Liceus. No Ceará, a denominação passou a ser Liceu Industrial do Ceará, nome que durou apenas um ano depois, em 1942, de acordo com o Decreto nº 4121, de 25 de fevereiro, recebeu o nome de Escola Industrial de Fortaleza.

A conjuntura nacional e internacional despertou o interesse do governo brasileiro em modernizar e melhorar o ensino profissional.

Em 1942, a Lei Orgânica do Ensino Industrial estabeleceu as bases da organização e do regime do ensino destinado à preparação profissional dos trabalhadores na indústria e definiu o ensino industrial como de 2º grau, em paralelo com o ensino secundário. Os cursos técnicos de três anos preparariam os alunos para uma nova modalidade de educação, que seria a formação técnica de segundo grau para a área industrial como atribuição das escolas técnicas industriais, que
naquele ano iniciaram suas atividades.

No estado do Ceará, a denominação Escola Técnica Federal do Ceará surge mediante a Lei nº 3.552 de 16 de fevereiro de 1953, alterada pelo Decreto-Lei nº 196, de 27 de agosto de 1969 vinculada ao MEC por intermédio da Secretaria de Educação Médio e Tecnológica - SEMTEC. É uma autarquia educacional, tendo se firmado no Estado como instituição de excelência no ensino técnico-profissional.

Cumpre salientar que tantas mudanças de nome foram decorrentes do sempre renovado papel da Instituição, para uma constante sintonia com os novos horizontes que eram delineados pela permanente dinâmica do progresso muito acelerada nas últimas décadas. A Escola Técnica Federal do Ceará teve inclusive seu campo de ação ampliado com a criação das UNED’s (Unidades Descentralizadas de Ensino) de Cedro e de Juazeiro do Norte (1994), viabilizando o ensino profissional em outras regiões do Estado.

A velocidade do desenvolvimento industrial do país e a inserção gradual de tecnologias avançadas demandam a formação de especialistas de diversos níveis, impondo um persistente reestudo na formação desses profissionais. Deste reestudo nascem os CEFET’s (Centros Federais de Educação Tecnológica) tendo por objetivo ministrar ensino em nível superior de graduação e pós-graduação, visando à formação de profissionais nas áreas de construção civil, industrial e tecnológica, a formação de professores e especialistas para o ensino médio e de formação profissional, formação de técnicos, promoção de cursos de extensão, aperfeiçoamento, atualização profissional e realização de pesquisas na área
técnico-industrial.

A denominação de Centro Federal de Educação Tecnológica do Ceará (CEFET-CE) foi oficializada pela Lei nº 8.948, de 8 de dezembro de 1994 e regulamentada pelo Decreto-Lei nº 2.406, de 27 de novembro de 1997 e pelo Decreto de 22/03/99 (DOU
de 22/03/99) que implantou a nova institucionalidade.

A necessidade de capacitação de novos profissionais levou o Governo Federal a sancionar a lei 11.892/08 que transformou os CEFET’s, Escolas Agrotécnicas e Técnicas em Institutos Federais de Educação, Ciência e Tecnologia (IF’s). Com o mesmo status das universidades federais, os IF’s serão obrigados a oferecer 20\% das vagas para a formação de professores, ou seja, os cursos de licenciaturas.

Os IF’s representam uma nova concepção da educação profissional e humana no Brasil e traduzem o compromisso do governo federal com os jovens e adultos. Esta nova rede de ensino tem um modelo institucional em que as unidades possuem autonomia administrativa e financeira. A nova instituição terá também forte inserção na área de Pesquisa e Extensão para estimular o desenvolvimento de soluções técnicas e tecnológicas.

O Instituto Federal de Educação, Ciência e Tecnologia do Ceará (IFCE) é uma Autarquia Educacional pertencente à Rede Federal de Ensino. Hoje, com 32 Campi, o Instituto Federal do Ceará se consolida como instituição de ensino inclusivo e de qualidade, cuja missão é produzir, disseminar e aplicar os conhecimentos científicos e tecnológicos na busca de participar integralmente da formação do cidadão, visando sua inserção social, política, cultural e ética.  O IFCE valoriza o compromisso ético com responsabilidade social, o respeito, a transparência e a excelência, em consonância com os preceitos básicos de cidadania e humanismo, com liberdade de expressão, cultura da inovação e idéias pautadas na sustentabilidade ambiental.

%\clearpage\section{BREVE HISTÓRICO DO CAMPUS MARACANAÚ}
\subsection{Breve Histórico do Campus Maracanaú}

Quando sancionada pelo Presidente Luis Inácio Lula da Silva, a Lei 11.892 criou trinta e oito Institutos Federais de Educação, Ciência e Tecnologia em 2008. O IFCE nasceu com nove Campi, dentre estes destaca-se o Campus Maracanaú, previamente estabelecido como CEFET - Uned Maracanaú em 2006.

Maracanaú é um município do Estado do Ceará que integra a região metropolitana de Fortaleza e constitui o maior Distrito Industrial do Ceará, caracterizada por um crescente contingente de empresas dos mais diversos setores, indústrias que vão desde o gênero alimentício e têxtil até a indústria metal-mecânica. O município de Maracanaú conta com aproximadamente 18.000 empresas instaladas (segundo dados do Governo do Estado do Ceará), ocupando a quarta posição no estado neste quesito.

O IFCE – Campus Maracanaú foi criado com o intuito de atender a demanda de mão de obra qualificada para as empresas do Estado do Ceará, favorecido por sua localização, estando mais próximo das indústrias em desenvolvimento e das já
existentes. Os cursos ofertados pelo campus Maracanaú sempre estiveram alinhados ao arranjo produtivo da região onde está inserido. Seu primeiro curso oferecido foi o curso técnico em Desenvolvimento de Software, ainda em 2006 como CEFET - Uned Maracanaú. Atualmente, o campus de Maracanaú do IFCE oferta ao todo nove cursos, distribuídos entre os níveis técnico e superior. Na pós-graduação, o campus possui o mestrado em Energias Renováveis e também atua no mestrado em Ciência da Computação, que funciona no campus de Fortaleza.

Neste contexto, o Eixo Tecnológico da Indústria, com apoio da Diretoria Geral do Campus Maracanaú, constituiu uma comissão regulamentada pelas portarias  \PORTARIACOMISSAO (ANEXO \ref{anexoPortariaComissao}) e  \PORTARIACOMISSAOALT (ANEXO \ref{anexoPortariaComissaoAlt}), para implantar 
%conforme resolução n° 020 de 13 de novembro de 2008, reproduzida no ANEXO \ref{anexoResolucao},
o \NOMECURSO, por considerar ser uma função estratégica dos setores industrial e de prestação de serviços.